\documentclass[11pt]{book}
\usepackage[paperwidth=145mm,paperheight=200mm,top=18mm,bottom=20mm,inner=20mm,outer=16mm]{geometry}
\usepackage{brumfiel}
% figures05.tex -- Chapter 5, "Measurement of Segments".
% All figures in this chapter are number-line / segment diagrams: labelled
% points on straight lines, with the unit segment shown below.  Every point
% position below is exact arithmetic on the unit u -- no point is placed by
% eye.  Proportions read off the iPad scans (scan05 pp.18-24, scan06 pp.6-7)
% and the photo PDF (PDF pp.87-97); see tools/constraints/ch05.py.
%
% Fig. 5-1 (a pictorial vignette of a man pacing off a building) is NOT
% redrawn -- see chapters/ch05-UNCERTAIN.md.
%
% Loaded here rather than in ch05.tex so that this file stands alone for
% tools/check_labels.py, whose driver inputs figures05.tex by itself.
\usetikzlibrary{decorations.pathreplacing}

% ---------------------------------------------------------------- Fig. 5-2
% Unit CD laid off from A: A1,A2,A3,A4 at 1..4 units; B between A3 and A4
% (measured at 0.36 of that gap, giving AB about 3.36 units).
\newcommand{\FIGVTWO}{%
\begin{bkfigure}[0.92\linewidth]
\begin{tikzpicture}[x=1.75cm,y=1cm,scale=1.257]
  \draw[given] (-0.20,0) -- (4.40,0);
  \foreach \x/\lb in {0/A,1/{A_1},2/{A_2},3/{A_3},3.36/B,4/{A_4}}{%
    \dt{\x,0}
    \node[vlab,above] at (\x,0) {$\lb$};}
  \node[vlab,right] at (4.40,0) {$l$};
  % unit segment CD, drawn below, spanning 0.52u to 1.52u.  Drop measured
  % from the book plate: gap between the two lines is 0.35 of the unit.
  \draw[given] (0.52,-0.61) -- (1.52,-0.61);
  \dt{0.52,-0.61} \dt{1.52,-0.61}
  \node[vlab,above,xshift=-2pt] at (0.52,-0.61) {$C$};
  \node[vlab,above,xshift=2pt] at (1.52,-0.61) {$D$};
  \node[slab,above] at (1.02,-0.63) {1};
\end{tikzpicture}
\figcap{5--2}
\end{bkfigure}}

% ---------------------------------------------------------------- Fig. 5-3
% The unit is too large: B lies between A and A1.  Re-measured off the book
% plate: B sits at 0.32 of AA1 (not 0.45), and the line runs on to 1.71 u.
\newcommand{\FIGVTHREE}{%
\begin{bkfigure}[0.62\linewidth]
\begin{tikzpicture}[x=1.75cm,y=1cm,scale=1.300]
  \draw[given] (-0.26,0) -- (1.71,0);
  \foreach \x/\lb in {0/A,0.32/B,1/{A_1}}{%
    \dt{\x,0}
    \node[vlab,above] at (\x,0) {$\lb$};}
\end{tikzpicture}
\figcap{5--3}
\end{bkfigure}}

% ----------------------------------------------------------- Figs. 5-4/5-5
% Archimedes' integer n = 4.  5-4: B = A4 exactly.  5-5: B between A3 and A4
% (measured at 0.18 of that gap).  In both plates the drawn line STARTS at A
% -- there is no stub to the left of A, unlike Fig. 5-2.
\newcommand{\FIGVFOURFIVE}{%
\begin{bkfigure}[0.94\linewidth]
\begin{tikzpicture}[x=1.75cm,y=1cm,scale=1.274]
  \draw[given] (0,0) -- (4.40,0);
  \foreach \x/\lb in {0/{A=A_0},1/{A_1},2/{A_2},3/{A_3},4/{A_4=B}}{%
    \dt{\x,0}
    \node[vlab,above] at (\x,0) {$\lb$};}
  \draw[given] (0.52,-0.61) -- (1.52,-0.61);
  \dt{0.52,-0.61} \dt{1.52,-0.61}
  \node[vlab,above,xshift=-2pt] at (0.52,-0.61) {$C$};
  \node[vlab,above,xshift=2pt] at (1.52,-0.61) {$D$};
\end{tikzpicture}
\figcap{5--4}
\end{bkfigure}%
\begin{bkfigure}[0.94\linewidth]
\begin{tikzpicture}[x=1.75cm,y=1cm,scale=1.274]
  \draw[given] (0,0) -- (4.40,0);
  \foreach \x in {0,1,2,3,3.18,4}{\dt{\x,0}}
  \foreach \x/\lb in {0/A,1/{A_1},2/{A_2},4/{A_4}}{%
    \node[vlab,above] at (\x,0) {$\lb$};}
  % A3 and B sit 0.18 u apart, as in the book; the labels are nudged apart
  % by 1.5pt each so they read as "A_3 B" rather than colliding.  The DOTS
  % stay at the exact positions 3 and 3.18 -- only the text moves.
  \node[vlab,above,xshift=-1.5pt] at (3,0) {$A_3$};
  \node[vlab,above,xshift=1.5pt] at (3.18,0) {$B$};
\end{tikzpicture}
\figcap{5--5}
\end{bkfigure}}

% ----------------------------------------------------------- Figs. 5-6/5-7
% 5-6: the unit interval CD divided into tenths (9 interior marks).
% 5-7: AB measured as about 2.3 units; the tenth-scale unit CD alongside.
\newcommand{\FIGVSIXSEVEN}{%
\begin{bkfigure}[0.90\linewidth]
\begin{tikzpicture}[x=7.5cm,y=1cm,scale=1.300]
  \draw[given] (0,0) -- (1,0);
  \dt{0,0} \dt{1,0}
  \node[vlab,above,xshift=-2pt] at (0,0) {$C$};
  \node[vlab,above,xshift=2pt] at (1,0) {$D$};
  \foreach \k in {1,...,9}{\draw[tick] ({\k/10},-0.085) -- ({\k/10},0.085);}
\end{tikzpicture}
\figcap{5--6}
\end{bkfigure}%
\begin{bkfigure}[0.90\linewidth]
\begin{tikzpicture}[x=1.9cm,y=1cm,scale=1.300]
  \draw[given] (0,0) -- (2.30,0);
  \foreach \x/\lb in {0/A,1/{A_1},2/{A_2},2.30/B}{%
    \dt{\x,0}
    \node[vlab,above] at (\x,0) {$\lb$};}
  % the unit CD alongside, itself divided into tenths
  \draw[given] (2.64,0) -- (3.64,0);
  \dt{2.64,0} \dt{3.64,0}
  \node[vlab,above,xshift=-2pt] at (2.64,0) {$C$};
  \node[vlab,above,xshift=2pt] at (3.64,0) {$D$};
  \foreach \k in {1,...,9}{\draw[tick] ({2.64+\k/10},-0.075) -- ({2.64+\k/10},0.075);}
\end{tikzpicture}
\figcap{5--7}
\end{bkfigure}}

% ---------------------------------------------------------------- Fig. 5-8
% Theorem 5-2: B'' chosen on the same side of A as B with AB'' = A'B'.
% B'' drawn between A and B (measured at 0.77 of AB).  AB is drawn the same
% length as A'B' (the theorem's given |AB| = |A'B'|), A sits directly under
% A', and the two rows are 0.17 of the segment length apart -- all measured
% off the book plate.
\newcommand{\FIGVEIGHT}{%
\begin{bkfigure}[0.72\linewidth]
\begin{tikzpicture}[x=4.6cm,y=1cm,scale=1.300]
  \draw[given] (0,0.80) -- (1,0.80);
  \dt{0,0.80} \dt{1,0.80}
  \node[vlab,above] at (0,0.80) {$A'$};
  \node[vlab,above] at (1,0.80) {$B'$};
  \draw[given] (0,0) -- (1,0);
  \foreach \x/\lb in {0/A,1/B}{%
    \dt{\x,0}
    \node[vlab,above] at (\x,0) {$\lb$};}
  \dt{0.77,0}
  \node[vlab,above,yshift=1.2pt] at (0.77,0) {$B''$};
\end{tikzpicture}
\figcap{5--8}
\end{bkfigure}}

% ---------------------------------------------------------------- Fig. 5-9
% Theorem 5-3: E on the same side of C as D with AB = CE; D between C and E
% (measured at 0.66 of CE).
\newcommand{\FIGVNINE}{%
\begin{bkfigure}[0.72\linewidth]
\begin{tikzpicture}[x=4.6cm,y=1cm,scale=1.300]
  \draw[given] (0,0.75) -- (1,0.75);
  \dt{0,0.75} \dt{1,0.75}
  \node[vlab,above] at (0,0.75) {$A$};
  \node[vlab,above] at (1,0.75) {$B$};
  \draw[given] (0,0) -- (1,0);
  \foreach \x/\lb in {0/C,0.66/D,1/E}{%
    \dt{\x,0}
    \node[vlab,above] at (\x,0) {$\lb$};}
\end{tikzpicture}
\figcap{5--9}
\end{bkfigure}}

% --------------------------------------------------------------- Fig. 5-10
% Change of scale: E is the midpoint of the unit CD, so AB = 3 with unit CD
% and AB = 6 with unit ED.  Horizontal layout re-measured off the book plate
% (all in CD units from C): A at 1.55, B at 4.55, the half-unit ED alongside
% at 5.15..5.65.  The earlier 2.24/5.24/5.92/6.42 spread the plate too wide.
\newcommand{\FIGVTEN}{%
\begin{bkfigure}[0.96\linewidth]
\begin{tikzpicture}[x=1.25cm,y=1cm,scale=1.300]
  % unit CD with midpoint E, braced and labelled 1
  \draw[given] (0,0) -- (1,0);
  \dt{0,0} \dt{1,0}
  \draw[tick] (0.5,-0.09) -- (0.5,0.09);
  \node[vlab,below] at (0,-0.02) {$C$};
  \node[vlab,below] at (0.5,-0.02) {$E$};
  \node[vlab,below] at (1,-0.02) {$D$};
  \draw[fig,line width=0.5pt,decorate,decoration={brace,amplitude=3pt}]
        (0,0.16) -- (1,0.16);
  \node[slab,above,yshift=2.6pt] at (0.5,0.19) {1};
  % AB, three units long; the "1" sits under its first third
  \draw[given] (1.55,0.55) -- (4.55,0.55);
  \dt{1.55,0.55} \dt{4.55,0.55}
  \node[vlab,above] at (1.55,0.55) {$A$};
  \node[vlab,above] at (4.55,0.55) {$B$};
  \draw[tick] (2.55,0.46) -- (2.55,0.64);
  \draw[tick] (3.55,0.46) -- (3.55,0.64);
  \node[slab,below,yshift=-2.2pt] at (2.05,0.53) {1};
  % the half-unit ED alongside
  \draw[given] (5.15,0.30) -- (5.65,0.30);
  \dt{5.15,0.30} \dt{5.65,0.30}
  \node[vlab,above] at (5.15,0.30) {$E$};
  \node[vlab,above] at (5.65,0.30) {$D$};
\end{tikzpicture}
\figcap{5--10}
\end{bkfigure}}


% ---- local macros (hoist into book preamble) ----
% Centred display headings ("REVIEW OF CHAPTER 5", "ALGEBRA REVIEW").
\newcommand{\Vhead}[1]{%
  \par\bigskip\begin{center}\bfseries\large\MakeUppercase{#1}\end{center}%
  \par\nopagebreak\medskip}
% Unnumbered star footnote, same reasoning as ch02's \IIstarnote: \footnote
% would set a numbered mark, but the book prints a bare "*". Used for the
% "irrational" footnote on book p.74.  NOTE: this chapter does NOT carry the
% "*Starred ... are optional" note -- book p.75, which holds this chapter's
% first starred exercises (Ex. Group 5-3, nos. 4 and 7), has no footnote at
% all; ch02 owns the book-wide first use.  Verified against PDF p.89.
\newcommand{\Vstarfoot}[1]{\begingroup\renewcommand{\thefootnote}{}%
  \footnotetext{*\,#1}\endgroup}
% Fig. 5-1 is a pictorial vignette, not a geometric diagram -> placeholder,
% per the plate rule (same treatment as Fig. 4-1). See ch05-UNCERTAIN.md.
\newcommand{\FIGVONE}{\begin{bkfigure}[0.44\linewidth]
\includegraphics[width=0.44\linewidth]{fig-5-1}%
\figcap{5--1}\end{bkfigure}}
% ---- end local macros ----

\begin{document}
\setcounter{chapter}{5}

\chapter*{\raggedright\Huge MEASUREMENT\\OF SEGMENTS}
\addcontentsline{toc}{chapter}{5.\ Measurement of Segments}
\markboth{MEASUREMENT OF SEGMENTS}{}
\vspace{-1em}\noindent\rule{\linewidth}{0.6pt}\vspace{1em}
\figurekey

%========================================================= 5-1
\section*{5--1\quad INTRODUCTION}
\addcontentsline{toc}{section}{5--1\quad Introduction}

In this chapter we are going to get a number for the length of a segment, that
is for the distance between two points. The way we get the length will be
exactly the same as the method we use in practical work. We will see that
finding a length is based upon simple \emph{counting}. For example, when you
pace off the length of a building (Fig.~5--1), you count how many steps you
take. Or when you say that the room is 15 feet wide, you mean that a foot ruler
can be laid down 15 times, one after the other, on a line. If the width of the
room is between 14 and 15 feet, the ruler can be laid down 14 times, but 15
times would put the end of the ruler past the wall. To measure the room more
accurately, you must choose a smaller unit of measure, say an inch, so that if
the width is 14 feet and three inches, the foot ruler can be laid down 14 times
and beyond that the inch ruler can be laid down three times. If even then you
have not come exactly to the wall of the room, you must choose an even smaller
unit of measure, some fraction of an inch, and so on.

\FIGVONE

\exercisegroup{5--1}

\begin{multicols}{2}
\begin{exlist}
\item Measure the length of your desk. Describe the process of measurement in
      terms of counting.
\item Measure the length of a page. Describe the counting you do to get the
      length.
\item Think of other examples of measurement and describe them in terms of
      counting.
\end{exlist}
\end{multicols}

%========================================================= 5-2
\section*{5--2\quad NUMBERS}
\addcontentsline{toc}{section}{5--2\quad Numbers}

Because we express length in terms of numbers, we need certain facts about
numbers that you learned in algebra and arithmetic. You have studied different
kinds of numbers. In the beginning grades you first used whole numbers
$1$, $2,\ldots$ and learned how to add, subtract, and multiply them. These
numbers we now call the \emph{positive integers} instead of ``whole numbers.''
When you tried to divide integers you found that it was necessary to have other
numbers, which you called \emph{fractions}, for example,
$\frac{1}{2}$, $\frac{2}{3}$, $\frac{5}{8}$, $\frac{7}{37}$. These are numbers
obtained by dividing one integer by another integer. \emph{The numbers which
are either positive integers or positive fractions are called positive rational
numbers.}

You also wrote numbers as \emph{decimals}, for example, 1.25, 2.537, 173.612.
Sometimes a rational number when written as a decimal ``comes out even,'' as
for example, $\frac{1}{2} = 0.5$, $\frac{3}{8} = 0.375$. Sometimes it does not
come out even, so that the decimal goes on forever, for example,
$\frac{1}{3} = 0.333\ldots$, $\frac{2}{3} = 0.6666\ldots$,
$\frac{1}{11} = 0.090909\ldots$. You also met with other numbers which did not
come out even, for example, $\sqrt{2} = 1.4142\ldots$,
$\sqrt{3} = 1.7321\ldots$. These numbers are different from $\frac{1}{3}$,
$\frac{2}{3}$, and $\frac{1}{11}$ because they not only do not come out even,
but \emph{the decimal never repeats}. In other words, no matter to how many
decimal places you were to calculate $\sqrt{2}$, you would never reach a place
where, from then on, one group of digits would occur over and over again in the
same order.

Decimals that, after some place, repeat one group of digits over and over are
called \emph{repeating decimals}. Decimals which come out even are called
\emph{terminating decimals}, and may be considered as repeating with zeros.

\exercise{5--2}

Write the following rational numbers as decimals to see that they are repeating
decimals. If they come out even (terminate), the repeating part is all zeros.
1/6, 5/6, 122/99, 1/13, 1/7, 6/5, 3/2.

\medskip
It is proved in advanced algebra that rational numbers are always repeating
decimals, and that, conversely, repeating decimals are rational numbers. (These
proofs are hard, so don't be discouraged if you cannot construct them; your
teacher may show you.) Hence, rational numbers and repeating decimals are the
same. What, then, about the decimals which never repeat? These numbers are not
rational numbers, and so they have been called \emph{irrational
numbers}.*%
\Vstarfoot{The word ``irrational'' is not a very pleasant choice
because it sounds like ``not sensible.'' What it really means is this. Rational
numbers are obtained by dividing one integer by another; that is, a rational
number is the ratio of two integers. Hence the word ``rational'' in ``rational
number'' means that the number is a \emph{ratio} of integers. Therefore, an
\emph{irrational} number is one which is \emph{not} the ratio of two integers.}
Irrational numbers are defined to be \emph{those numbers which are not
rational}. Examples of irrational numbers are $\sqrt{2}$, $\sqrt{3}$,
$\sqrt{5}$, $\sqrt[3]{2}$, $\sqrt[3]{17}$, and, in fact, any root of an integer
which is not itself an integer.

All the decimals together are defined to be the \emph{real numbers}. Thus a
real number is a decimal. If the decimal repeats, the number is a rational
number. If the decimal does not repeat, it is an irrational number.

\exercisegroup{5--3}
% Book p.75 holds this chapter's first starred exercises (nos. 4 and 7) but
% prints NO "*Starred ... are optional" footnote -- verified against PDF p.89.
% ch02 owns the book-wide first use; deliberately no \starnote here.

\begin{multicols}{2}
\begin{exlist}
\item Write the symbols for several rational numbers and for several irrational
      numbers.
\item Is $3.131313\ldots$, where the 1's and 3's keep alternating, a rational
      number?
\item Is $0.234234234\ldots$, where 2, 3, and 4 keep following each other, a
      rational number?
\item[\stex 4.] The number $1.010010001\ldots$ is formed from just 0's and 1's.
      After each 1 there is one more 0 than the previous time. Is this number
      rational?
\item[5.] Give the definitions of rational number, irrational number, real
      number.
\item[6.] A real number is either rational or irrational. Why?
\item[\stex 7.] There is a rule for deciding when a fraction written as a
      decimal will have its repeating part all zeros and when it will not. For
      example, each of the fractions 3/20, 17/250, 5/16, and 4/25 has its
      repeating decimal part all zeros. The fractions 7/15, 5/12, and 11/23 do
      not. Can you discover the rule?
\end{exlist}
\end{multicols}

When we write a decimal, for example 1.23, we mean $1 + 0.2 + 0.03$ or
$1 + 2/10 + 3/100$. Similarly, $4.383 = 4 + 3/10 + 8/100 + 3/1000$. Thus, any
terminating decimal (one which comes out even) may be written as a sum of
ordinary fractions with denominators 10, $100,\ldots$. If the decimal is
repeating, this sum is an \emph{infinite} one, for example,
$1/3 = 0.333\ldots = 3/10 + 3/100 + 3/1000 + 3/10000 + \cdots$.

\exercisegroup{5--4}

\begin{multicols}{2}
\begin{exlist}
\item Write the following decimals as sums of rational numbers: 4.82, 1.29,
      2.75, 7.122, 0.8875, 2.40. In order to express each number as a single
      fraction, add the rational numbers.
\item Write the following numbers as decimals: $3 + 2/10 + 4/100$,
      $2 + 3/10 + 5/100 + 6/1000$, $4 + 3/100$, $1 + 1/10 + 1/10000$.
\item Write the following repeating decimals as sums of rational numbers:
      $0.6666\ldots$, $0.7777\ldots$, $0.121212\ldots$.
\item Write the following numbers as decimals:
      $1 + 2/10 + 2/100 + 2/1000 + \cdots$,
      $3 + 1/10 + 4/100 + 1/1000 + 4/10000 + \cdots$,
      $8/10 + 8/100 + 8/1000 + 8/10000 + \cdots$.
\end{exlist}
\end{multicols}

%========================================================= 5-3
\section*{5--3\quad INEQUALITIES}
\addcontentsline{toc}{section}{5--3\quad Inequalities}

In this book you are expected to know how to do arithmetic, that is, how to
add, subtract, multiply, and divide real numbers. But there is one part of
arithmetic you have not yet studied systematically. That part is concerned with
comparing the sizes of real numbers. Since you already know what it means to
say that the number $x$ is larger than the number $y$, all we need is a symbol
to indicate this. We therefore make the definition: $x > y$ means ``$x$ is
greater than $y$.'' Another way to say this is ``$x - y$ is a positive
number.'' We also define $x < y$ to mean ``$x$ is less than $y$,'' or, in other
words ``$y$ is greater than $x$.'' Listed below are several properties of real
numbers. These properties can be established as indicated in the review of
Chapter 4.

\begin{enumerate}[label=(\alph*), leftmargin=2.2em, itemsep=1pt, topsep=3pt]
\item If $x > y$ and $y = z$, then $x > z$.
\item If $x > y$ and $y > z$, then $x > z$.
\item $x$ is a positive number if and only if $x > 0$.
\item $x$ is a negative number if and only if $x < 0$.
\item If $x > y$, then $x - y > 0$, and conversely.
\end{enumerate}

\medskip
We define $x \geqq y$ to mean ``either $x > y$ or $x = y$,'' and we read
``$\geqq$'' as ``is greater than or equal to.'' We make a similar definition
for $x \leqq y$ (``$x$ is less than or equal to $y$'').

\begin{enumerate}[label=(\alph*), leftmargin=2.2em, itemsep=1pt, topsep=3pt,
                  start=6]
\item If $x > y$ and $u \geqq v$, then $x + u > y + v$.
\end{enumerate}

\exercisegroup{5--5}

\begin{multicols}{2}
\begin{exlist}
\item Substitute numbers for the variables in properties (a), (b), (e), and (f)
      above, and thus verify their truth for certain cases.
\item Write in symbols: $a$ is greater than $b$; $a$ is greater than or equal
      to $b$; $u$ is less than $w$; $a + c$ is less than or equal to $r + s$;
      $(a + b)^2$ is positive; $z$ is a negative number.
\item If $x$ is any real number, $x^2 \geqq 0$. Why?
\item Why is it true that if $x > 0$ and $y > 0$, then $xy > 0$?
\item Tell why $xy < 0$ if $x > 0$ and $y < 0$.
\end{exlist}
\end{multicols}

%========================================================= 5-4
\section*{5--4\quad MEASUREMENT}
\addcontentsline{toc}{section}{5--4\quad Measurement}

To measure a line segment we must first tell what our unit of length is. Some
units of length often used in practical work are inch, foot, yard, mile, meter,
centimeter, etc. Having picked a unit of length, we will describe the geometric
process whereby we get a real number, that is, a decimal for the length of a
line segment.

Choose any segment $CD$ as a unit (Fig.~5--2). This segment and all segments
congruent to it are defined to have length 1. Suppose $AB$ is any segment on a
line $l$. On the same side of $A$ as $B$ there is a point $A_1$ such that
$AA_1 \cong CD$. If $A_1 = B$, then the length of $AB$ is 1.

\FIGVTWO

\FIGVTHREE

\noindent If $B$ lies between $A$ and $A_1$ (Fig.~5--3), then our unit was too
large to measure the length of $AB$ in integers (whole numbers). We will return
to this possibility later. If $A_1$ lies between $A$ and $B$, then there is a
point $A_2$, on the same side of $A_1$ as $B$, such that either $A_2 = B$, or
$B$ is between $A_1$ and $A_2$, or $A_2$ is between $A_1$ and $B$. If
$A_2 = B$, the length of $AB$ is 2. If $B$ is between $A_1$ and $A_2$, then our
unit will not measure the length in integers. If $A_2$ is between $A_1$ and
$B$, there is a point $A_3$, on the same side of $A_2$ as $B$, such that
$A_2A_3 \cong CD$.

The process described above is what we actually follow when we measure. If we
lay out segments congruent to $CD$ from $A$ toward $B$, we know from experience
that we eventually either reach $B$ exactly, so that the length of $AB$ is an
integer, or reach a point \emph{beyond} $B$. Thus in Fig.~5--2, $A_4$ is beyond
$B$ but $A_3$ is still between $A$ and $B$. The length in this case will be a
number between 3 and 4. If the above process can always be carried out in our
geometry, then we have described the beginning of the process of measurement.

The only part of the above process that does not follow from our postulates is
the proof that we can always get to or beyond $B$. Because we need to do this,
we now add a new postulate to our list. It was the Greek mathematician
Archimedes (287--212 B.C.) who first stated this postulate. He apparently did
not see how to get it from Euclid's postulates and, needing it, assumed it was
true. It is one of the things that Euclid forgot to state as a postulate. It is
called \emph{Archimedes' axiom} or \emph{Archimedes' postulate} in spite of its
use by Eudoxus much earlier. This is the first of the two new postulates in
Postulate Group IV which we need for measurement. The other one will be given
later.

\medskip
IV--1 \emph{Archimedes' postulate.} If $AB$ is any segment on a line $l$ and
$CD$ any other segment, then there is a definite number $n$ of points
$A_1$, $A_2,\ldots$, $A_n$ on $l$ such that the points $A = A_0$, $A_1$,
$A_2,\ldots$, $A_n$ are arranged in this order, and all the segments $AA_1$,
$A_1A_2,\ldots$, $A_{n-1}A_n$ are congruent to $CD$, and either $B = A_n$ or
$B$ lies between $A_{n-1}$ and $A_n$.

\medskip
The integer $n$ is called \emph{Archimedes' integer}. When $n$ is 4, an
illustration would look like Fig.~5--4 or Fig.~5--5.

\FIGVFOURFIVE

\exercisegroup{5--6}

\begin{multicols}{2}
\begin{exlist}
\item If $AB$ is this segment \raisebox{0.18ex}{$A$\rule[0.35ex]{2.2em}{0.4pt}$B$}
      and $CD$ is this segment
      \raisebox{0.18ex}{$C$\rule[0.35ex]{1.3em}{0.4pt}$D$}, find how many
      points are needed to get beyond $B$.
\item Draw a segment $AB$ and a segment $CD$. Find the integer $n$ of
      Archimedes' axiom for your segments.
\end{exlist}
\end{multicols}

Using Archimedes' postulate, we find that the length of $AB$ is either an
integer or is between two integers. If the length lies between two integers,
the next step in the process is to divide the unit interval into tenths
(Fig.~5--6). (We would not have to keep to tenths if we were not trying to come
out with a decimal.) Then we count ``tenths'' until either we reach $B$ exactly
or $B$ lies between two of the tenth marks. If we reach $B$ exactly, then the
length of $AB$ comes out exactly in tenths. If $B$ is not reached exactly, then
we must divide the unit interval into one-hundredths and repeat the process. In
this way we obtain a real number or decimal, for the length of $AB$. In
Fig.~5--7 the length of $AB$ is about 2.3.

\FIGVSIXSEVEN

We have indicated a segment by writing the two endpoints of the segment as
$AB$. At this point we also need a notation for the \emph{length} of the
segment. For this we will use $\sg{AB}$, that is $\sg{AB} = $ \emph{the length
of the segment} $AB$. Whenever you see the symbol $\sg{AB}$ you must remember
that it represents a \emph{number}, while the symbol $AB$ stands for a \emph{set
of points}.

Except for dividing a segment into tenths, the entire procedure of this section
can be carried out (that is, described) in our geometry. At the moment we do
not have the means for dividing a segment into tenths. As soon as we have
studied parallel lines (Chapter 8), we will be able to make this construction.
Since we do not need the results of this chapter until after we have studied
parallel lines, our procedure is logically correct.

\exercisegroup{5--7}

\begin{multicols}{2}
\begin{exlist}
\item Draw a line segment. Then draw a segment that will be your unit. Now
      measure the first segment, describing how you use Archimedes' postulate
      and relations of betweenness in the process.
\item Draw a line segment. Measure it with your ruler to the nearest 1/16 of an
      inch. Describe your measurement process, using Archimedes' postulate and
      betweenness.
\end{exlist}
\end{multicols}

%========================================================= 5-5
\section*{5--5\quad POINTS ON A LINE}
\addcontentsline{toc}{section}{5--5\quad Points on a Line}

We have just described the measurement process for getting a number for the
length of a segment. We now turn the question around and ask whether, given a
number, there is a segment with that number for its length. Of course, we would
like this to be true, but our postulates so far do not tell us that it is true.
The process of measurement describes a method for ``sneaking up'' on the point
$B$, but if we start out with a number we are not assured that there is a point
to sneak up on! To have this property, we must postulate it. This postulate
tells us that there are enough points on the line.

\medskip
IV--2. \emph{The completeness postulate.} For every line $l$ and any point $A$
on $l$, and for any positive real number $x$, there is, on a given side of $A$,
a point $B$ such that $\sg{AB} = x$.

\medskip
The name of the last postulate is meant to signify that our lines are complete,
that is, that they have all the points on them we expect them to have.

\exercisegroup{5--8}

\begin{multicols}{2}
\begin{exlist}
\item Draw a segment that is to be the unit segment. Draw another segment that
      is 1.3 units in length.
\item With the unit segment of Exercise 1, draw a segment of length
      $2 + \frac{1}{2} + \frac{1}{4} + \frac{1}{8}$ units.
\end{exlist}
\end{multicols}

%========================================================= 5-6
\section*{5--6\quad LENGTH AND CONGRUENCE}
\addcontentsline{toc}{section}{5--6\quad Length and Congruence}

When we started talking about congruent segments we had in mind a notion of
size. Now that we have lengths for segments it should turn out that congruent
segments have the same length, and conversely that segments of the same length
are congruent. We now prove these theorems. To do so we need the following
theorem (4--3) that we proved in Chapter 4. \emph{If $AB \cong A'B'$ and $ACB$
and if $C'$ is on the same side of $A'$ as $B'$ and $AC \cong A'C'$, then
$A'C'B'$.}

\medskip\noindent
\textbf{*Theorem 5--1.}\ \emph{If $AB \cong A'B'$, then
$\sg{AB} = \sg{A'B'}$.}
\refstepcounter{theorem}
\par\medskip

\noindent\emph{Proof.} To get the length of segment $AB$, we establish a
sequence of points \emph{between} $A$ and $B$. Let us designate this sequence
by $C_1$, $C_2$, $C_3,\ldots$. Now, consider the corresponding sequence,
$C_1'$, $C_2'$, $C_3',\ldots$, where $AC_1 \cong A'C_1'$,
$AC_2 \cong A'C_2'$, $AC_3 \cong A'C_3',\ldots$, and all the points $C_1'$,
$C_2'$, $C_3',\ldots$ are on the same side of $A'$ as $B'$. Using Theorem
4--3, we see that all the points $C_1'$, $C_2'$, $C_3',\ldots$, are between
$A'$ and $B'$; therefore, this sequence determines the same real number
(decimal) as does the sequence $C_1$, $C_2$, $C_3,\ldots$.

\medskip\noindent
\textbf{*Theorem 5--2.}\ \emph{If $\sg{AB} = \sg{A'B'}$, then
$AB \cong A'B'$.}
\refstepcounter{theorem}
\par\medskip

\noindent\emph{Proof.} This theorem is the converse of Theorem 5--1. We will
prove it by choosing a point $B''$, on the same side of $A$ as $B$, such that
$AB'' \cong A'B'$, and then proving that $B'' = B$ (Fig.~5--8). This will show
that $AB \cong A'B'$. Either $B'' = B$ or $B'' \neq B$. Since
$AB'' \cong A'B'$, then from Theorem 5--1 $\sg{AB''} = \sg{A'B'}$. We are given
that $\sg{AB} = \sg{A'B'}$. Hence $\sg{AB''} = \sg{AB}$. Thus we must show that
if $B'' \neq B$, we cannot have $\sg{AB''} = \sg{A'B'}$. But if $B'' \neq B$,
then, if the unit segment is smaller than $BB''$, the corresponding Archimedes
integers for $AB$ and $AB''$ would be different. This would lead to different
lengths for $AB$ and $AB''$. Hence $B''$ must be the same as $B$ and
$AB \cong A'B'$.

\FIGVEIGHT

These two theorems give us the right to use the words ``the same length'' in
place of ``congruent.'' Therefore we can use everyday language in talking about
congruence of segments.

\exercisegroup{5--9}

\begin{multicols}{2}
\begin{exlist}
\item $AB \cong CD$ and $CD \cong EF$, and $\sg{AB} = 7.26$ inches. What is
      $\sg{EF}$?
\item If $\sg{PQ}$ is 1.23 ft, and $VU \cong RS$, and $\sg{RS} = 1.23$ ft, is
      $QP \cong VU$?
\item If $AB \cong CD$, $CD \cong EF$, $EF \cong GH$, is $\sg{AB} = \sg{GH}$?
\item An isosceles triangle can be defined now in two ways. State the two
      definitions.
\end{exlist}
\end{multicols}

%========================================================= 5-7
\section*{5--7\quad MORE ABOUT INEQUALITIES}
\addcontentsline{toc}{section}{5--7\quad More about Inequalities}

In Chapter 4 we defined $>$ and $<$ for segments. In this chapter (Section
5--3) we defined $>$ and $<$ for numbers. Since ``same length'' means
``congruent,'' it should also be true that if $AB > CD$, then
$\sg{AB} > \sg{CD}$, and conversely.

\medskip\noindent
\textbf{*Theorem 5--3.}\ \emph{If $AB > CD$ then $\sg{AB} > \sg{CD}$, and
conversely, if $\sg{AB} > \sg{CD}$ then $AB > CD$.}
\refstepcounter{theorem}
\par\medskip

\noindent\emph{Proof.} The theorem could also be stated as ``$AB > CD$ if and
only if $\sg{AB} > \sg{CD}$.'' Or, again, as ``If $AB > CD$ then
$\sg{AB} > \sg{CD}$, and conversely.'' We first prove that $AB > CD$ implies
$\sg{AB} > \sg{CD}$. The hypothesis, $AB > CD$, means that if $E$ is on the
same side of $C$ as $D$, and $AB \cong CE$, then $D$ is between $C$ and $E$.
(See Fig.~5--9.) By Theorem 5--1, $\sg{AB} = \sg{CE}$. In the measurement
process, if the unit is chosen smaller than $DE$, the corresponding Archimedes'
integers for $CD$ and $CE$ will be different, and in fact the integer for $CD$
will be smaller than that for $CE$. Hence $\sg{CD} < \sg{CE}$ and, since
$\sg{CE} = \sg{AB}$, $\sg{CD} < \sg{AB}$.

\FIGVNINE

The converse is proved in a similar manner. We have given that
$\sg{AB} > \sg{CD}$. Choose $E$, on the same side of $C$ as $D$, such that
$AB \cong CE$. By Theorem 5--1, $\sg{AB} = \sg{CE}$. Now either $CDE$,
$D = E$, or $CED$. (Why?) We cannot have $D = E$. (Why?) An argument like the
first part of the proof will show we cannot have $CED$. Hence we must have
$CDE$. It follows from the definition of ``less than'' for segments that
$CD < CE \cong AB$.

\exercisegroup{5--10}

\begin{multicols}{2}
\begin{exlist}
\item If $ABC$ and $\sg{AB} = 5$, what can you say about $\sg{AC}$?
\item If $AB > EF$ and $\sg{EF} = 4$, what can you say about $\sg{AB}$?
\item If $AB > CD > EF$, what can you say about $\sg{AB}$ and $\sg{EF}$?
\item If $ABCD$, what can you say about $\sg{AD}$ and $\sg{BC}$?
\item If $\sg{AB} > \sg{AC}$, and $C$ is a point on the line containing $A$ and
      $B$, and $C$ lies on the same side of $A$ as $B$, where is the point $C$?
\end{exlist}
\end{multicols}

%========================================================= 5-8
\section*{5--8\quad CHANGE OF SCALE}
\addcontentsline{toc}{section}{5--8\quad Change of Scale}

In measurement we choose a fixed segment for a unit of length. What happens if
a different segment is chosen for the unit? This is a familiar situation
because we are all used to dealing with inches, feet, yards, or miles to
measure length. In Fig.~5--10, if $CD$ is the unit, the length of $AB$ is 3. If
$ED$ is the unit, the length of $AB$ is 6. This illustrates the following
theorem.

\FIGVTEN

\medskip\noindent
\textbf{Theorem 5--4.}\ \emph{If $\sg{AB} = x$ when $CD$ is the unit, and if
$\sg{CD} = y$ when $EF$ is the unit, then $AB = xy$ when $EF$ is the unit.}
\refstepcounter{theorem}
\par\medskip

We will not give the proof of this theorem for arbitrary $x$ and $y$ because it
is long to write out and quite hard if $y$ is an irrational number. But if $y$
is an integer, the proof is very easy and is really clear from Fig.~5--10,
where $y = 2$.

\exercisegroup{5--11}

\begin{multicols}{2}
\begin{exlist}
\item $\sg{AB} = 5$ when $CD$ is the unit, and $\sg{CD} = 2.3$ when $EF$ is the
      unit. What is $\sg{AB}$ when $EF$ is the unit?
\item The width of the room is 14 ft. What is the width in inches? How does
      this follow from Theorem 5--4?
\item One kilometer is 1000 meters. One meter is 100 centimeters. How many
      centimeters in a kilometer? Show how this follows from Theorem 5--4.
\item $\sg{PQ} = 6.5$ when $\sg{RS} = 1$. $\sg{RS} = \frac{1}{2}$ when
      $\sg{UV} = 1$. What is $\sg{PQ}$ when $\sg{UV} = 1$?
\item If $CD < EF$, how does $\sg{AB}$ with unit $CD$ compare with $\sg{AB}$
      with unit $EF$? In other words, if we use two different units and compute
      two lengths for one segment, does the smaller unit give the smaller or
      the larger length?
\end{exlist}
\end{multicols}

%========================================================= Review
\Vhead{Review of Chapter 5}
\addcontentsline{toc}{section}{Review of Chapter 5}
\markboth{REVIEW OF CHAPTER 5}{}

\begin{exlist}
\item In this chapter we have given definitions of the following: real number,
      rational number, irrational number, positive integer, greater than, less
      than, the process of measurement. Be sure you know these definitions.
\item This chapter has two new postulates, Archimedes' postulate and the
      completeness postulate. Go over them again and be sure you see what they
      mean.
\item In this chapter we have proved the following theorems. Study them again
      to be sure you know what they say. You should be able to state clearly
      the facts upon which each proof depends.
\end{exlist}

\begin{align*}
(AB \cong CD) &\iffx \sg{AB} = \sg{CD},\\
(AB > CD) &\iffx \sg{AB} > \sg{CD},\\
(AB < CD) &\iffx \sg{AB} < \sg{CD}.
\end{align*}

\begin{note}
In this chapter we have shown that congruence is equivalent to equality of
length. From this point on in our development of geometry you may use the
ordinary language of length in speaking of congruence.
\end{note}

%========================================================= Algebra review
\Vhead{Algebra Review}
\addcontentsline{toc}{section}{Algebra Review}

\begin{exlist}
\item Segment $XY$ has a length of 28 units. Point $T$ is between $X$ and $Y$
      and $\sg{XT} = \sg{YT} + 9$. Determine the lengths of $XT$ and $YT$.
\item Points $R$, $T$, and $S$ are in this order on a line. Segment $RT$ is
      twice as long as segment $TS$. The difference in their lengths is 11.
      Find the length of $RS$.
\item Points $A$, $B$, and $C$ are on one line. The length of $AC$ is 12 more
      than that of $AB$. The length of $BC$ is four times the length of $AB$.
      Find the lengths of $AB$ and $AC$ if (a) $A$ is between $B$ and $C$, and
      (b) $B$ is between $A$ and $C$.
\item Points $A$, $B$, $C$, and $D$ are in this order on one line.
      $\sg{BC} = \sg{AB} + 4$; $\sg{CD} = \sg{AB} + \sg{BC}$; $\sg{AD} = 32$.
      Find $\sg{AB}$.
\item The length of $XY$ is $\frac{2}{3}$ of the length of $XZ$. The length of
      $XZ$ is $\frac{3}{5}$ of the length of $YZ$. If these three points are
      collinear: (a) What is the order of the points on $l$? (b) Express
      $\sg{XY}$ as a fraction of $\sg{YZ}$. (c) Determine $\sg{XY}$ if
      $\sg{XZ} - \sg{XY} = 6$.
\item The point $M$ is the midpoint of segment $AB$, $N$ is the midpoint of
      $MB$, and $O$ is the midpoint of $NB$. If $\sg{AN} - \sg{NO} = 20$, find
      the length of $AB$.
\item $\sg{AC} = 2\sg{AB}$, $\sg{AB} + \sg{BC} = \sg{AC} + 2$. If
      $\sg{AB} > 1$, prove that the three points $A$, $B$, and $C$ are not on
      one line.
\item A square and equilateral triangle have equal perimeters. The length of a
      side of the triangle exceeds the length of a side of the square by six
      inches. Find the perimeter of each.
\item The sum of the perimeters of a square and a rectangle is 36 ft. The
      rectangle is twice as long as it is wide, and the width of the rectangle
      is 1 ft more than the length of a side of the square. Find the dimensions
      of each.
\item The point $B$ is the midpoint of segment $AC$. Point $X$ is between $B$
      and $C$, and $\sg{BX} = 3$ inches. Show that a square having a side
      congruent to $AB$ has an area of 9 in$^2$ more than a rectangle with
      sides congruent to $AX$ and $XC$.
\end{exlist}

\end{document}
